\documentclass[11pt]{article}
\usepackage{fontspec}
\setmainfont{Times New Roman}
\setsansfont{Arial}
\setmonofont{Consolas}
\usepackage{amsmath,amssymb}
\usepackage{geometry}
\usepackage{microtype}
\geometry{a4paper,margin=3cm}
\setlength{\parindent}{1.25em}
\setlength{\parskip}{0pt}
\emergencystretch=4em
\allowdisplaybreaks
\let\A\relax
\newcommand{\A}{\mathfrak A}
\newcommand{\R}{\mathfrak R}

\begin{document}

\begin{center}
\textbf{\S\ 14. \quad Formy devętogo poręda (Sistem III).}
\end{center}
\[
\text{Svijanje }\left\{\begin{array}{l}
\theta\cdot K \text{ s }H,\\
K,j,\A \text{ s }L,\\
j,\A \text{ s }\sigma,\\
f \text{ s }H^2.
\end{array}\right.
\]

\noindent\textit{Nereducibilne formy:}
\[
\begin{array}{c@{\quad}l@{\qquad}c@{\quad}l}
\mathrm{A)}&(\vartheta\cdot k,\eta,x)^4,\quad
H_k(\vartheta\cdot k,\eta,x)^4&
\mathrm{B)}&\begin{array}{l}
(KLu)^3,\quad K_i(KLu)^3,\quad K_i^2,\\
(klx)^3,\quad K_i(Klx)^3
\end{array}\\[0.7em]
\mathrm{C)}&L_j,\quad L_j^2&
\mathrm{D)}&\A_i,\\[0.7em]
\mathrm{E)}&(j\sigma x)&
\mathrm{G)}&(aH^2u)^3,\quad (aH^2u)^4,\quad a_\eta(aH^2u)^3.
\end{array}
\]

\noindent\textit{Redukcije:}

\noindent\mbox{A)} Formy $H_\vartheta(k\eta x)^3$, $H_\vartheta^2(k\eta x)^2$, $H_\vartheta^2(k\eta x)^3$ razkladajut se kako nastale črez poslědovateljno jednorazovo svijanje s razložimymi formami.

\noindent\mbox{B)} Rędy form $K_iL_k$, $K_i^2(KLu)$, $K_i^2(klx)$:

reducenty: $\theta_\eta H_\vartheta$, $a_\eta^2(aHu)$, $\theta_\eta^2(\vartheta\eta x)$.

Formy $L_i$, $K_i(KLu)$, $K_i(klx)$, $(KLu)^2$, $(klx)^2$ razkladajut se kako transvekcije nad funkcionalnymi determinantami (\S\ 3).

Formy $L_k$, $L_k(KLu)$, $L_k(KLu)^2$, $L_k(klx)$, $L_k(klx)^2$, $L_k^2$, $L_k^2(KLu)$, $L_k^2(klx)$, $L_k^3$, $K_i(KLu)^2$, $K_i(klx)^2$ razkladajut se kako nastale črez jednorazovo svijanje s razložimymi formami:
\[
\begin{gathered}
L_\vartheta,\quad H_k(KHu),\quad L_\vartheta(\vartheta lx),\quad H_k^2,\\
L_\vartheta^2,\quad L_\vartheta^2(\theta Lu),\quad
K_\eta(KHu),\quad \theta_i(\vartheta lx),
\end{gathered}
\]
po \S\ 3. 2), \mbox{a)} i \mbox{b)}.

\noindent\mbox{C)} Ręd form $(jlx)^3$: reducent $(j\eta x)^3$.

Formy $(jlx)^2$ i $L_j(jlx)$ nastale sut črez jednorazovo svijanje s razložimoju formoju $(j\eta x)^2$.

\noindent\mbox{D)} Ręd form $\A_i(\A Lu)$: reducent $\A_\nu^2$.

Formy $(\A Lu)^2$, $(\A Lu)^3$: jednorazovo svijanje s razložimoju formoju $(\A Hu)^2$.

\noindent\mbox{E)} Ręd form $(j\sigma x)^2$: reducent $a_\sigma^2$ črez zaměnu $j$ nižšimi formami rěda form (\S\ 5):
\[
a_\sigma^2(a\theta u)^2=\frac13(j\sigma x)^2,\qquad
a_\sigma^2(a\theta\sigma x)^2=\frac13(j\sigma x)^4.
\]

\noindent\mbox{F)} Ręd form $\A_\sigma^2$: reducent $a_\sigma^2$.

Forma $\A_\sigma$: reducent $a_\sigma^2$ črez zaměnu $\A$ nižšimi formami rěda form:
\[
a_\sigma a_\R^2=\frac23\A_\sigma\pmod{\nu}.
\]

\noindent\textit{Poslědki:} Jedino forma $j$ vhodi v svijanje s $\sigma$ i $(\nu\sigma x)$.

$N$ ne vhodi v svijanje s $L$, bo forma $N_i$, odpovědajuča formě $\A_i$, razkladaje se.

\end{document}
